\oldpage[I]{88}
Enfin, si $M$ est somme directe de deux $A$-modules $N$, $P$, il est immédiat
que $\widetilde{M}=\widetilde{N}\oplus\widetilde{P}$~; toute somme directe
étant limite inductive de sommes directes finies, les assertions de (ii) sont
démontrées.

On notera que (\hyperref[I.1.3.8-fr]{1.3.8}) prouve que les faisceaux
isomorphes aux faisceaux associés aux $A$-modules forment une
\emph{catégorie abélienne} (T, I, 1.4).

On notera aussi qu'il résulte de (\hyperref[I.1.3.9-fr]{1.3.9}) que si $M$ est
un $A$-module \emph{de type fini}, c'est-à-dire s'il existe un homomorphisme
surjectif $A^n\to M$, alors il existe un homomorphisme surjectif
$\widetilde{A^n}\to\widetilde{M}$, autrement dit le $\widetilde{A}$-Module
$\widetilde{M}$ est \emph{engendré par une famille finie de sections au-dessus
de $X$} (\hyperref[0.5.1.1-fr]{0, 5.1.1}), et réciproquement.

\begin{env}[1.3.10]
\label{I.1.3.10-fr}
Si $N$ est un sous-module d'un $A$-module $M$, l'injection canonique
$j:N\to M$ donne par (\hyperref[I.1.3.9-fr]{1.3.9}) un homomorphisme injectif
$\widetilde{N}\to\widetilde{M}$, qui permet d'identifier canoniquement
$\widetilde{N}$ à un \emph{sous-$\widetilde{A}$-Module} de $\widetilde{M}$~;
nous supposerons toujours faite cette identification. Si $N$ et $P$ sont deux
sous-modules de $M$, on a alors
\begin{equation}
  \widetilde{(N+P)}=\widetilde{N}+\widetilde{P}
  \tag{1.3.10.1}\label{I.1.3.10.1-fr}
\end{equation}
\begin{equation}
  \widetilde{(N\cap P)}=\widetilde{N}\cap\widetilde{P}
  \tag{1.3.10.2}\label{I.1.3.10.2-fr}
\end{equation}
car $N+P$ et $N\cap P$ sont respectivement l'image de l'homomorphisme
canonique $N\oplus P\to M$, et le noyau de l'homomorphisme canonique
$M\to(M/N)\oplus(M/P)$ et il suffit d'appliquer
(\hyperref[I.1.3.9-fr]{1.3.9}).

On conclut de (\hyperref[I.1.3.10.1-fr]{1.3.10.1}) et
(\hyperref[I.1.3.10.2-fr]{1.3.10.2}) que si
$\widetilde{N}=\widetilde{P}$, on a $N=P$.
\end{env}

\begin{corollary}[1.3.11]
\label{I.1.3.11-fr}
Dans la catégorie des faisceaux isomorphes aux faisceaux associés aux
$A$-modules, le foncteur $\Gamma$ est exact.
\end{corollary}

En effet, soit
$\widetilde{M}\xrightarrow{\widetilde{u}}\widetilde{N}
\xrightarrow{\widetilde{v}}\widetilde{P}$ une suite exacte correspondant à
deux homomorphismes $u:M\to N$, $v:N\to P$ de $A$-modules. Si
$Q=\operatorname{Im}u$ et $R=\operatorname{Ker}v$, on a
$\widetilde{Q}=\operatorname{Im}\widetilde{u}
=\operatorname{Ker}\widetilde{v}=\widetilde{R}$
(cor. (\hyperref[I.1.3.9-fr]{1.3.9})), donc $Q=R$.

\begin{corollary}[1.3.12]
\label{I.1.3.12-fr}
Soient $M$, $N$ deux $A$-modules.
\begin{enumerate}
  \item[\rm(i)] Le faisceau associé à $M\otimes_A N$ s'identifie
  canoniquement à $\widetilde{M}\otimes_{\widetilde{A}}\widetilde{N}$.
  \item[\rm(ii)] Si de plus $M$ admet une présentation finie, le faisceau
  associé à $\operatorname{Hom}_A(M,N)$ s'identifie canoniquement à
  $\mathscr{H}\!om_{\widetilde{A}}(\widetilde{M},\widetilde{N})$.
\end{enumerate}
\end{corollary}

(i) Le faisceau
$\mathscr{F}=\widetilde{M}\otimes_{\widetilde{A}}\widetilde{N}$ est associé
au préfaisceau
\[
  U\mapsto\mathscr{F}(U)=\Gamma(U,\widetilde{M})
  \otimes_{\Gamma(U,\widetilde{A})}\Gamma(U,\widetilde{N})
\]
$U$ parcourant la base (\hyperref[I.1.1.10-fr]{1.1.10, (i)}) de $X$ formée
des $D(f)$, où $f\in A$. Or, $\mathscr{F}(D(f))$ s'identifie canoniquement à
$M_f\otimes_{A_f}N_f$ en vertu de (\hyperref[I.1.3.7-fr]{1.3.7}) et
(\hyperref[I.1.3.6-fr]{1.3.6}). On sait par ailleurs que le $A_f$-module
$M_f\otimes_{A_f}N_f$ est canoniquement isomorphe à $(M\otimes_A N)_f$
(\hyperref[0.1.3.4-fr]{0, 1.3.4}), qui lui-même est canoniquement isomorphe à
$\Gamma(D(f),\widetilde{(M\otimes_A N)})$
(\hyperref[I.1.3.7-fr]{1.3.7} et \hyperref[I.1.3.6-fr]{1.3.6}). En outre, on
vérifie aussitôt que les isomorphismes canoniques
\[
  \mathscr{F}(D(f))\simeq\Gamma(D(f),\widetilde{(M\otimes_A N)})
\]
